% Transcription — fonds Alexandre Grothendieck, shelfmark 94, batch 1 (pages 1-20).
% Established from the facsimile archives/batches/94/batch-01.pdf (a copy of
% 94.pdf fetched through the project's relay,
% https://grothendieck-relay.onrender.com/source/94.pdf; PDF page k =
% archivists' page k, no cover sheet), per the skill transcribe-grothendieck.
% The folder is twenty-three pages; this is the first of two batches (pages
% 21-23 are batch 2, transcribed separately; it was read before this file was
% finished and its notation is followed: f_*, V(F), \underline{\mathrm{Hom}},
% \underline{\mathrm{Ext}}, statement numbers running on to 18-21).
% Pass: Opus 5.5 (claude-opus-5-5), 2026-09-26 — first pass, unchecked against the pages by a human.
%
% Pages skipped: 11 and 13, typed leaves unrelated to the notes and with
% nothing in his hand (a typescript chapter on valued fields, and a typed
% report on a redaction). Page 1 is a brown cover inscribed in ink, top right;
% it takes no \page marker, its words become the section heading, with a \note.
%
% Pages summarised: none. No typed leaf carries his ink.
%
% Pages transcribed from his hand: 2-10, 12, 14-20. Blue ink, fast drafting;
% pp. 10 and 12 on darker (tan) paper. Formulas and numbered statements carry;
% part of the connecting prose does not (pp. 6, 7, 17 margin, 19 end are the
% worst).
%
% His pagination: none on the leaves. His numbering of statements runs
% through the batch and on into batch 2: Lemmes 1-2, Prop. 3-4, Lemme 5,
% Remarque 5 bis, Lemme 6, "Cor. Main" 7, Lemme 8, "Cor. Main" 9, 10, a
% cancelled 11 (p. 12) and a new 11 (p. 14), Théorème 12, Remarque 13,
% Théorème 14, Lemmes 15, 16, 17; batch 2 carries 18-21.
%
% What the batch is about. One run, not two. The cover (p. 1) reads
% « Construction des Modules cohérents par leurs foncteurs représentatifs.
% Application à la représentabilité d'un foncteur ... », and p. 2 is headed
% « Construction de faisceaux de modules », the second half of the folder
% title. Nothing on these pages is headed « Dualité des modules »; the duality
% material (a « faisceau fondamental » Omega on a Cohen-Macaulay scheme,
% injective hulls, pp. 8-14) sits inside the same numbered run as a tool, so
% the title's first part may name another leaf or the archivists' reading of
% this one. Despite the inventory group « Géométrie et topologie combinatoire »,
% this is algebraic geometry: representability of functors
% F : Mod(A) -> (Ens) by finitely generated modules. Lemma 1: descent along a
% faithfully flat B. Lemma 2 / Prop. 3: A noetherian, criteria (left
% exactness, commutation with filtered colimits, completion condition, F_f and
% F_{A/f^n} representable); Prop. 4 for A local via injective hulls; Lemma 5:
% pro-representability with stationary localisations; Remarque 5 bis:
% constructibility of p -> rank F(k(p)) instead. Lemma 6, Cor. 7, Lemmas 8-11:
% generic Tor-vanishing against regular sequences, and Hom into a dualising
% Omega commuting with reduction mod a system of parameters, used to show a
% pro-system is stationary. Théorème 12 (A noetherian, quotient of a regular
% ring): F representable by a module of finite type iff (ii) left exact, (i)
% filtered colimits, (iii) a/b completion conditions, (iv) finiteness, (v)
% constructibility, (vi) base change of F(Omega) mod parameter systems.
% Théorème 14: for f : X -> S proper of finite presentation, F, G of finite
% presentation, G S-flat, there is P of finite presentation on S with
% Hom(P, M) = f_* Hom(F, G (x) M), compatible with base change; proof by
% th. 12, with Lemmas 15-17 (generic base change for Ext^i, Tor, R^j f_*) —
% the P that batch 2's Théorème 19 uses.
%
% Readings to check first: the cover's second paragraph (p. 1); the end of
% p. 5 and the proof on p. 6; the middle of p. 7; the letter read Q on p. 9
% (probably Omega); the marginal NB of p. 17; the cancelled « condition (v) »
% block on p. 17; the end of p. 19. Batch 2's references to « th. 12 » and to
% its condition (ii) as left exactness agree with p. 15.
%
% Apparatus count: 85 \ill{}, 49 \uncertain{}.

\documentclass[11pt,a4paper]{article}
% Le préambule commun à toutes les transcriptions du fonds Grothendieck.
%
% Il tient en une idée : **l'incertitude se compose, elle ne se résout pas.**
% Une transcription qui lisse un mot douteux a détruit la seule chose pour
% laquelle on la faisait — un lecteur doit pouvoir savoir, à chaque endroit,
% ce qui était sur le feuillet et ce qui a été ajouté. Les six macros qui
% suivent sont donc l'appareil critique, et non de la décoration.
%
% Les mêmes commandes sont reconnues par `scripts/render.mjs`, qui produit la
% vue de lecture du site. Une macro ajoutée ici doit l'être aussi là-bas,
% faute de quoi le rendu échouera bruyamment — ce qui est voulu : un rendu
% silencieusement faux est pire qu'un rendu absent.

\ProvidesPackage{grothendieck}[2026/08/08 Transcriptions du fonds Grothendieck]

\usepackage{amsmath,amssymb,amsthm}
\usepackage{mathrsfs}
\usepackage{stmaryrd}
% Les diagrammes commutatifs sont partout dans ce fonds : c'est la langue
% même de la géométrie algébrique de Grothendieck, et une transcription qui
% les rendrait en prose ne transcrirait rien.
\usepackage{tikz-cd}
% Français par défaut — c'est la langue des feuillets — mais l'anglais est
% chargé aussi : la traduction et le résumé sont des documents anglais, et la
% typographie française (espaces avant les deux-points) y serait une faute.
% Ces documents basculent par \selectlanguage{english} en tête de corps.
\usepackage[english,french]{babel}
\usepackage{fontspec}
\usepackage{geometry}
\geometry{a4paper,margin=25mm,marginparwidth=28mm}
\usepackage{marginnote}
\usepackage{xcolor}
% ulem, not soul. `soul`'s \st and \ul are built on a character-by-character
% scan that XeLaTeX's font machinery breaks: the document fails with an
% undefined control sequence rather than degrading. ulem does the same two jobs
% and is Unicode-engine safe. `normalem` stops it hijacking \emph, which the
% transcriptions use for Grothendieck's own emphasis.
\usepackage[normalem]{ulem}
\usepackage{hyperref}
\hypersetup{colorlinks=true,linkcolor=black,urlcolor=black,pdfborder={0 0 0}}

% --- Métadonnées du lot -----------------------------------------------------
% Reprises telles quelles par le script de rendu, qui les affiche en tête de
% la vue de lecture. Elles fixent ce que le fichier prétend couvrir : sans
% elles, une transcription flotte sans point d'attache dans un fonds de seize
% mille feuillets.

\newcommand{\folder}[1]{\gdef\grothFolder{#1}}
\newcommand{\batch}[1]{\gdef\grothBatch{#1}}
\newcommand{\leaves}[2]{\gdef\grothFirst{#1}\gdef\grothLast{#2}}
\newcommand{\foldertitle}[1]{\gdef\grothFoldertitle{#1}}
\gdef\grothFolder{?}\gdef\grothBatch{?}\gdef\grothFirst{?}\gdef\grothLast{?}\gdef\grothFoldertitle{}

% --- Le résumé --------------------------------------------------------------
%
% Il ouvre la lecture modernisée, sous le titre « Résumé », et il est composé
% en petit. Ce n'est pas un ornement : c'est le seul passage du document qui
% ne soit pas des mathématiques, et un lecteur doit pouvoir le reconnaître
% avant de l'avoir lu — puis le sauter s'il connaît déjà le sujet, ou s'y
% arrêter s'il ne le connaît pas. Le filet qui le suit marque la reprise du
% corps.

\newenvironment{resume}
  {\small\setlength{\parskip}{0.45em}}
  {\par\medskip\hrule\medskip}

% --- Mots-clés ---------------------------------------------------------------
%
% En anglais, en clôture du résumé de la lecture modernisée : le vocabulaire
% moderne sous lequel chercher ce que les feuillets construisent. Le manifeste
% du site (`npm run manifest`) les extrait pour étiqueter la cote entière —
% cette ligne est la source unique des tags, il n'y a pas de fichier à côté.
\newcommand{\keywords}[1]{\par\smallskip\noindent
  {\footnotesize\textsc{Keywords} --- \textit{#1}}\par}

% --- L'appareil critique ----------------------------------------------------

% Le numéro de feuillet, en marge. C'est l'ancre qui permet de régler un
% désaccord sur une formule en regardant *un* feuillet plutôt qu'une chemise
% de six cents. Le site s'en sert aussi pour tourner les pages du fac-similé
% quand on fait défiler la transcription.
\newcommand{\leaf}[1]{%
  \marginnote{\footnotesize\color{gray!70}\textsf{#1}}%
  \label{leaf:#1}\ignorespaces}

% Illisible : on ne devine pas. Un mot inventé qui se lit comme les autres est
% le pire résultat possible.
\newcommand{\ill}{\textcolor{red!60!black}{[\,\ldots\,]}}

% Lecture douteuse : le mot est donné, et signalé comme incertain.
\newcommand{\uncertain}[1]{\uline{#1}}

% Ajout de l'éditeur — une lettre manquante, un mot sous-entendu.
\newcommand{\add}[1]{\textcolor{blue!50!black}{[#1]}}

% Biffé par Grothendieck lui-même. Ce qu'il a rayé fait partie du document :
% une rature dit souvent où la pensée a changé de direction.
\newcommand{\struck}[1]{\sout{#1}}

% Note du transcripteur, sur l'état matériel du feuillet.
\newcommand{\note}[1]{\par\smallskip{\small\color{gray!60!black}\textit{[#1]}}\par\smallskip}

% Note marginale de Grothendieck, distincte du corps du texte.
\newcommand{\marginal}[1]{\par\smallskip\noindent\hspace{1em}%
  {\small\color{gray!60!black}#1}\par\smallskip}

% --- Titre ------------------------------------------------------------------

\AtBeginDocument{%
  \begin{center}
    {\large\bfseries Fonds Alexandre Grothendieck --- cote n\textsuperscript{o}~\grothFolder}\\[2pt]
    {\itshape\grothFoldertitle}\\[4pt]
    {\small feuillets \grothFirst--\grothLast\ (lot \grothBatch)}\\[2pt]
    {\footnotesize\color{gray!60!black}Transcription établie d'après le fac-similé numérisé
     par l'Université de Montpellier.\\
     Les lectures incertaines sont soulignées, les illisibles notées [\,\ldots\,],
     les ajouts entre crochets.}
  \end{center}
  \vspace{1em}}

\endinput


\folder{94}
\batch{1}
\pages{1}{20}
\dating{[à partir de 1965]}
\watermark{Édition de démonstration}
\foldertitle{Dualité des modules / Construction de faisceaux de modules : notes manuscrites (s.d.)}

\begin{document}

\section*{Construction des Modules cohérents par leurs foncteurs représentatifs}
\note{titre repris de la couverture (page 1), de sa main, à l'encre, dans
l'angle supérieur droit d'un feuillet brun par ailleurs vierge : « Construction
des Modules cohérents \struck{et de schémas finis sur d'autres par descente}
par leurs foncteurs représentatifs », puis, sous un trait : « Application à la
représentabilité d'un foncteur $G/S$ \ill{} \ill{} d'une section, à la
définition des invariants $\Omega^{1}(y)$\,\ldots ». Le $G$ est suivi d'une
lettre biffée ; « $\Omega^{1}(y)$ » est une lecture incertaine.}

\page{2}
\note{la page porte en tête, soulignée, « Construction de faisceaux de
modules » — la seconde moitié du titre du dossier.}

\subsection*{Construction de faisceaux de modules}

$A$ anneau \struck{noeth.}

$F \colon \mathrm{Mod}(A) \to (\mathrm{Ens})$ foncteur.

Pour toute algèbre $B$ sur $A$, on en conclut un foncteur
\[
  F_B \colon \mathrm{Mod}(B) \to (\mathrm{Ens})
\]
(composé de la restriction $\mathrm{Mod}(B) \to \mathrm{Mod}(A)$ et de $F$).
Si $F$ repr.\ par $P$, $F_B$ l'est par $P_B$.
\note{le composé est dessiné : une flèche oblique « restr. » de
$\mathrm{Mod}(B)$ vers $\mathrm{Mod}(A)$, puis une flèche montante vers
$(\mathrm{Ens})$.}

Si $f \in A$, on pose $F_f = $ \uncertain{le} \uncertain{nom} de $F_{A_f}$.

\emph{Lemme 1.} \struck{Supposons $F$ exact à gauche, et soit} Soit $B$ une
$A$-algèbre fid.\ plate. Alors
\[
  F \text{ représentable (resp.\ par un $A$-module de type fini, de prés.\ finie)}
\]
\[
  \Updownarrow
\]
\[
  F \text{ exact à gauche, et } F_B \text{ représentable (resp.\ par un $B$-module de type fini, de prés.\ finie).}
\]
\note{« Soit $B$ une $A$-algèbre fid.\ plate » est écrit au-dessus de la
ligne biffée ; « $A$- » et « $B$- » devant « module » sont ajoutés en
interligne.}

\emph{Lemme 2.} Soit $f \in A$, supposons $A$ noeth. Alors $F$ repr.\
\struck{\ill{}} par \struck{un} $A$-module de t.f.\ ssi il satisfait :
\begin{itemize}
  \item[(i)] $F$ exact à gauche ;
  \item[(ii)] $F(\widehat{M}) \xrightarrow{\ \sim\ } \varprojlim F(M/f^{n}M)$
    si $M$ un $A$-module de t.f.\ \struck{($\widehat{M} = \varprojlim M/f^{n}M$)}
    ($\widehat{A} = \varprojlim A/f^{n}A$) ;
  \item[(iii)] $F$ commute aux $\varinjlim$ filtrantes ;
  \item[(iv)] $F_f$ repr.\ \struck{par $A$} ;
  \item[(v)] $\forall n$, $F_{A/f^{n}A}$ représentable.
\end{itemize}
\note{« Lemme 2 » est écrit au-dessus d'un « Corollaire 2 » biffé. Dans (ii),
le chapeau est sur le $M$ de $F(\widehat{M})$ ; la parenthèse biffée et la
précision « de t.f. » sont en interligne, la définition de $\widehat{A}$
ajoutée à droite et reliée par un trait.}

Nécessité claire, (iii) exprimant le fait que \struck{le} module $P$ qui
représente $F$ est de t.f.\ (donc de prés.\ finie).

\page{3}
\emph{Suffisance.} \uncertain{Grâce} à (ii) et Lemme 1 on est ramené à voir que
$F_{\widehat{A}}$ est représentable, car posant $B = \widehat{A} \times A_f$,
$F_B$ sera alors représentable en vertu de (iv), donc $F$ l'est ; le $P$ le
représentant sera alors de t.f.\ grâce à (iii).
\note{les numéros de conditions de cette page sont entourés ; le premier est
lu (ii) sans certitude.}

Pour tout $n$, $F_{A/f^{n+1}A}$ est représentable, par (v), par un module
$P_n$ sur $A_n$ \struck{\ill{}} de type fini grâce à (iii). Soit
$\widehat{P} = \varprojlim P_n$, qui est de t.f.\ sur $\widehat{A}$ (N.B.\ les
$P_n$ « se recollent »). \struck{\ill{}} Utilisant (ii), on trouve que
$\widehat{P}$ représente $F_{\widehat{A}}$ sur la catégorie des $A$-modules de
t.f., et utilisant (iii), on trouve que $\widehat{P}$ représente
$F_{\widehat{A}}$ partout.
\note{« (v) », entouré, est ajouté au bout de la ligne ; $A_n$, en interligne
sous l'indice $A/f^{n+1}A$ qu'il souligne, abrège cet anneau.}

\emph{Proposition 3.} $A$ noeth. Pour que $F$ soit repr.\ par un module de
t.f., il faut et il suffit que $F$ satisfasse aux conditions suivantes :
\begin{itemize}
  \item[(i)] $F$ exact à g. ;
  \item[(ii)] $F$ commute aux $\varinjlim$ \uncertain{monomorphes} filtrantes ;
  \item[(iv)] Pour tout quotient $B$ de $A$, \struck{et tout $f \in B$,} il
    existe un $f \in B$, non nilpotent, tel que $F_{B_f}$ soit représentable ;
  \item[(iii)] Pour tout $f \in A$, et tout module de type fini $M$ sur
    $\widehat{A}$ ($\widehat{A} = \varprojlim A/f^{n+1}A$),
    \[
      F(M) \xrightarrow{\ \sim\ } \varprojlim F(M/f^{n+1}M)
    \]
    \struck{où $\widehat{M} = \varprojlim M/f^{n+1}M$.}
\end{itemize}
\note{la condition numérotée (iv), d'abord écrite (iii), précède sur la page
la condition (iii) ; une flèche en marge ramène (iii) au-dessus d'elle. Dans
(iii), le $F(M)$ est écrit en surcharge d'un $F(\widehat{M})$. Le mot sous
« $\varinjlim$ » en (ii) est ajouté en interligne.}

\page{4}
Nécessité claire. Pour la suffisance, on note que les conditions sont stables
par passage à une algèbre quotient $B$. Par récurrence noeth., on peut donc
supposer que pour tout tel quotient $B = A/J$, avec $J \neq 0$, $F_B$ est
représentable (par un module de t.f.). On peut supposer aussi $A \neq 0$. Il
existe par (iv) un $f \in A$ non nilpotent, tel que $F_f$ soit repr. D'autre
part, l'hyp.\ de récurrence noeth., prenant pour tout $n$, où
$f^{n+1} \neq 0$, montre que $F_{A/f^{n+1}A}$ est repr. On peut donc appliquer
le Lemme 2, qui donne la conclusion.
\note{le « (iv) » est entouré ; l'indice de $F_f$ est repassé.}

En pratique, la condition (iv) est parfois difficile à vérifier. On va donner
d'autres conditions.
\note{un double trait horizontal sépare ce paragraphe de ce qui précède.}

\emph{Proposition 4.} $A$ local noeth. Pour que $F$ soit repr.\ par module de
t.f., il faut et il suffit que $F$ satisfasse les conditions
\begin{itemize}
  \item[(i)] $F$ exact à gauche ;
  \item[(ii)] $F$ commute aux $\varinjlim$ ;
  \item[(iii)] Pour tout $\widehat{A}$-module de t.f.\ $M$,
    \[
      F(\widehat{M}) \xrightarrow{\ \sim\ } \varprojlim F(M/\mathfrak{m}^{n+1}M),
    \]
    et
  \item[(iv)] $F(\widehat{M})$ est de t.f.\ sur $\widehat{A}$ (il suffit
    \ill{} si $M$ est \uncertain{fini} sur $A$).
\end{itemize}
Nécessité claire. Pour suffisance, grâce au Lemme 1 et condition (ii), on
peut supposer $A$ complet.
\note{dans (iii) et (iv), le chapeau de $\widehat{M}$ est ajouté après coup ;
un premier mot de (iv) est biffé.}

\page{5}
En vertu de (ii), sa restriction aux modules finis sur $A_n = A/\mathfrak{m}^{n+1}$
\struck{\ill{}} est pro-représentable, \struck{par $(P_n)$, $P_n$ module sur
$A_n$,} et en vertu de la condition (iv), testée avec l'injectif
indécomposable $I_n$ de $A_n$, on trouve qu'il est représentable,
\struck{\ill{} par $P_n$,} par un module $P_n$ de t.f.\ sur $A_n$. Les $P_n$
se recollent, soit $P = \varprojlim P_n$, qui est un $A$-module de t.f. On
conclut de (iii) que $P$ représente $F$ sur les modules de t.f., et en vertu
de (ii) sur tous les modules, cqfd.
\note{« finis sur $A_n = A/\mathfrak{m}^{n+1}$ » est ajouté au-dessus de la
ligne, en remplacement d'un mot biffé ; l'indice de $\mathfrak{m}$ est lu sans
certitude. Le « (ii) » de la dernière phrase est repassé et pourrait être un
autre numéro.}

\emph{Lemme 5.} $A$ noeth. \struck{Conditions équivalentes.}
Supposons qu'il \uncertain{satisfait} les conditions
\begin{itemize}
  \item[(i)] $F$ exact à g. ;
  \item[(ii)] $F$ commute aux $\varinjlim$ filtrantes ;
  \item[(iii)] $\forall\, \mathfrak{p} \in \operatorname{Spec} A$, et tout $M$
    de t.f.\ \uncertain{sur} $\widehat{A}_{\mathfrak{p}}$, on a
    \[
      F(M) \xrightarrow{\ \sim\ } \varprojlim F(M/\mathfrak{m}_{\mathfrak{p}}^{n+1}M) ;
    \]
  \item[(iv)] $\forall$ module de \struck{type fini} longueur finie $M$ sur
    $A_{\mathfrak{p}}$, $F(M)$ est un $A_{\mathfrak{p}}$-module de t.f. ;
  \item[(v)] $\forall$ idéal premier $\mathfrak{p}$, $\exists\, f \in
    A/\mathfrak{p}$ \uncertain{non nul} tel que $F_{(A/\mathfrak{p})_f}$ soit
    représentable.
\end{itemize}
Alors $F$ pro-représentable par un syst.\ projectif strict $(P_i)$ de modules
de t.f., tel que pour tout $\mathfrak{p} \in \operatorname{Spec} A$, le syst.\
des $(P_i)_{\mathfrak{p}}$ soit stationnaire (i.e.\ $F_{\mathfrak{p}}$ est
représentable par un module de t.f.\ sur $A_{\mathfrak{p}}$).
\note{sur la page, la conclusion (« $F$ pro-représentable \ldots ») est écrite
d'abord, sous « Lemme 5 » (qui surcharge « Proposition », biffé), et un long
trait courbe la renvoie après « Alors », qui suit (v) ; la phrase introduisant
les conditions (« Supposons qu'il \ldots\ les conditions ») est écrite au-dessus
d'un « 2) \ill{} » biffé. Le « $(P_i)$ » est ajouté au-dessus de la ligne ;
la parenthèse « (i.e.\ \ldots) » est dans la marge de droite. En (i), le mot
« exact » est repassé. En (v), les « $A/\mathfrak{p}$ » sont lus sur des
« $A_{\mathfrak{p}}$ » surchargés.}

\emph{Dém.} Pour tout $\mathfrak{p} \in S = \operatorname{Spec} A$, la prop.~4
indique que $F_{\mathfrak{p}}$ est représentable par $P_{\mathfrak{p}}$ de
t.f.\ sur $A_{\mathfrak{p}}$ (grâce à (iii) à (iv)). Les $P_{\mathfrak{p}}$
sont \ill{} par

\page{6}
généralisation. Pour \uncertain{prouver} la \uncertain{pro}-repr., il suffit
de prouver (\uncertain{compte} tenu de (i) et (ii)) que pour tout
$\alpha \in F(M)$, $M$ de t.f.\ sur $A$, il existe un « \uncertain{module} de
définition $N$ », i.e.\ un plus petit sous-module $N$ de $M$ tel que
$\alpha \in \operatorname{Im}(F(N) \to F(M))$. Or on a pour tout
$\mathfrak{p} \in S$ un $N_{\mathfrak{p}} \subset M_{\mathfrak{p}}$
($= \operatorname{Im}(\alpha_{\mathfrak{p}} \colon P_{\mathfrak{p}} \to
M_{\mathfrak{p}})$). \struck{Soit $U$} Il suffit de vérifier que le système des
$N_{\mathfrak{p}}$ provient d'un sous-module $N$ de $M$ (on sait alors aisément
qu'il convient). Soit $U \subset S$ \uncertain{le} plus grand \ill{} sur
lequel se \uncertain{recolle}, soit $\underline{N}_U \subset \widetilde{M}|U$ ;
\struck{\ill{}} si on avait $U \neq S$, soit $x \in S - U$ un pt maximal de
$S - U$. Soit $S' = \operatorname{Spec} \mathcal{O}_{S,x}$,
\struck{Il existe un voisinage \ill{}} $U' = U \cap S'$, \struck{alors}
$N_{S'} = \widetilde{N}_x$ (N.B.\ $x = \mathfrak{p}$ !), alors
$\widetilde{N}_x | U' = \underline{N}_U | U'$, donc $\exists$ voisinage
\uncertain{ouvert} $V$ de $x$, \struck{\ill{}} et un
$\underline{N}' \subset \widetilde{M}|V$, tel que
$\underline{N}'_x = N_x \subset M_x$ et
$\underline{N}' | V \cap U = \underline{N}_U | V$, i.e.
\[
  \begin{cases}
    \underline{N}'_x = N_x \subset M_x \\
    \underline{N}'_s = N_s \subset M_s & \text{si } s \in U \cap V.
  \end{cases}
\]
\note{« et (ii) », au-dessus de la première ligne, est ajouté en interligne ;
les deux numéros sont entourés. « $\underline{N}_U \subset \widetilde{M}|U$ »
est ajouté au-dessus de la ligne ; le mot biffé sous « on avait » est lu
\uncertain{complètement}. « (N.B.\ $x = \mathfrak{p}$ !) » : le
« $x = \mathfrak{p}$ » est entouré. La première ligne du système,
« $\underline{N}'_x = N_x \subset M_x$ », est aussi écrite, un peu plus haut,
avec un point sur le $M$.}

Prenant $V$ ($= \operatorname{Spec} A_f$) assez petit, on peut supposer que
\struck{la condition (v)} $F_{(A/\mathfrak{p})_f}$ est représentable
(\uncertain{grâce} à la condition (v)), \struck{\ill{} Il s'ensuit que}
\uncertain{et} on veut tout de suite \ill{} (quitte à \ill{} encore $V$) que,
$\alpha_V \in F(\underline{N}')$ donc \ill{} $N_s \subset \underline{N}'_s$
pour tout $s \in V$ ; \uncertain{or} \ill{} $\underline{N}'_s$ \ill{}
l'image de $N_s$ dans $M_s$ \struck{est égal} est tout \uncertain{entier}
$\ldots\ \mathfrak{p}$, donc est tout pour $s \in V(\mathfrak{p})$ par
Nakayama. Cela \uncertain{suffit}.
\marginal{$\longrightarrow$ Donc on est ramené au cas $\underline{N}' = M$
\struck{\ill{}}}
\note{la fin de la page est très remaniée : « $= \operatorname{Spec} A_f$ »
est au-dessus du $V$, « (v) \ill{} » en interligne relié à « la condition » ;
à droite, un cadre porte « $\alpha_V \in F(\underline{N}')$ donc
$N_s \subset \underline{N}'_s$ \ill{} pour tout $s \in V$ ». L'ordre de
lecture des fragments est incertain ; la note marginale, à gauche, est reliée
au passage par une flèche.}

\page{7}
\emph{Remarque 5 bis.} Avec les notations du lemme~5, et en présence des
conditions (i) à (iv), la condition (v) équivaut à la condition suivante, qui
se prête mieux à une vérification directe :
\begin{itemize}
  \item[(v bis)] La fonction $\mathfrak{p} \mapsto
    \operatorname{rang}_{k(\mathfrak{p})} F(k(\mathfrak{p}))$ sur
    $\operatorname{Spec} A$ est constructible.
\end{itemize}
Il faut cependant supposer que $A$ satisfait à la condition : pour tout
quotient intègre $B$ de $A$, $\mathrm{Nor}(B)$ contient un ouvert non vide.
\note{cette dernière phrase est ajoutée en interligne, serrée, et reliée par un
trait courbe ; « $\mathrm{Nor}(B)$ » est une lecture incertaine (lieu normal ?).}

Le fait que (v) $\Rightarrow$ (v bis) est trivial.

En sens inverse, on utilise le fait que $F$ définit un « syst.\ local » de
modules $P_{\mathfrak{p}}$. Dans le cas $A$ intègre, \struck{\ill{}} si le
rang des $P_{\mathfrak{p}} \otimes k(\mathfrak{p})$ est constant et égal à $n$
(dans un ouvert non vide), \struck{\ill{}} on peut pour
$\mathfrak{p} \in U$ \uncertain{\ill{}} \ill{} les $P_{\mathfrak{p}}$
($\mathfrak{p} \in U$) sont libres de rang $n$. D'ailleurs, \struck{\ill{}}
comme $F$ commute aux $\varinjlim$, \uncertain{on voit} qu'on peut \ill{} un
\ill{} $\subset N$, \uncertain{avec} \ill{} que ces $P_{\mathfrak{p}}$ soient
\uncertain{tous} $\subset N$, avec $N$ \ill{} \ill{} $M$. Utilisant le fait
que \struck{\ill{}} $F(N)$ est de type fini, on trouve aisément que pour tous
les $\mathfrak{p} \in U$ de \uncertain{codimension} 1, sauf un nombre fini, on
a $P_{\mathfrak{p}} = N_{\mathfrak{p}}$. Quitte à restreindre $U$, on peut
supposer que ceci est vrai pour tous les $\mathfrak{p}$ de codim.\ 1 de $U$.
Mais alors, si on prend $U$ \uncertain{normal}, on conclut que
$P_{\mathfrak{p}} = N_{\mathfrak{p}}$ pour \emph{tout} $\mathfrak{p} \in U$,
\struck{donc} d'où facilement que $N$ représente $F|U$.
\note{la page est d'une lecture difficile au milieu : l'ordre des phrases
autour de « D'ailleurs » est incertain. « sauf un nombre fini » est souligné
d'un trait qui le relie à la ligne suivante.}

\page{8}
\emph{Lemme 6.} Soient $S$ schéma loc.\ noeth., $F$ \struck{\ill{}} coh.\ sur $S$.

a) Soit $J$ le nilradical de $\mathcal{O}_S$ ($S_0 = S_{\mathrm{red}}$). Il
existe un ouvert dense $U$ de $S$ tel que \struck{\ill{}}
$\mathrm{gr}_J(\mathcal{O}_S)$ et $\mathrm{gr}_J(F)$ loc.\ libres sur $U_0$.

b) Soit $U$ comme dessus. Alors pour tout $s \in U$, et \struck{toute} suite
$\mathcal{O}_{S,s}$-régulière $f_1 \ldots f_n$ de $\mathfrak{m}_s$,
\struck{posant $M = \mathcal{O}_{S,s}/(f_1, \ldots, f_n)$}, on a
\[
  \mathrm{Tor}_i^{\mathcal{O}_{S,s}}\bigl(F_s,\ \mathcal{O}_{S,s}/(f_1, \ldots, f_n)\bigr) = 0
  \quad \text{pour } i > 0 .
\]
\note{« loc. » (devant « noeth. ») surcharge un mot biffé ;
« $S_0 = S_{\mathrm{red}}$ » et « et $\mathrm{gr}_J(F)$ » sont ajoutés en
interligne, ce dernier en remplacement d'une rature illisible. Dans le Tor,
l'indice de $F_s$ surcharge un autre symbole.}

\emph{Dém.} L'assertion a) est triviale. Pour b), on est ramené au cas où
$S = U$, et où $S = \operatorname{Spec} \mathcal{O}_{S,s}$, \struck{et $F$
provient d'un faisceau loc.\ libre sur $S_0$}. Pour calculer les
$\mathrm{Tor}_i$ on utilise le complexe de chaînes de Koszul
$K_{\cdot}(\underline{f})$, qui est une résolution de $A/(f_1, \ldots, f_n)$
($A = \mathcal{O}_{S,s}$). On est ramené à prouver que les images
$f_1 \ldots f_n$ des $f_i$ dans $A_0$ forment \struck{aussi} une suite
$M$-régulière. En fait, on prouve un peu plus que, posant
$N = M/(f_1, \ldots, f_n)M$, on a $\mathrm{gr}_J(N)$ \struck{loc.\ libre}
libre. C'est bien standard si $n = 1$, et on continue par récurrence.
\note{« et où $S = \operatorname{Spec} \mathcal{O}_{S,s}$ » est ajouté en
interligne. Le $M$ de « $M$-régulière » n'est pas défini ailleurs que dans la
ligne biffée de l'énoncé ; il faut sans doute y lire le module $F_s$ ou son
image.}

\emph{\uncertain{Cor. Main} 7.} $S$ schéma noethérien, $P$ et $\Omega$
faisceaux cohérents. Alors $\exists$ ouvert dense $U$ de $S$ tel que pour tout
$s \in U$ et toute \uncertain{suite} $\mathcal{O}_{S,s}$-régulière
$f_1 \ldots f_n$,
\note{le titre, souligné, est lu comme au dossier~90 (« CorMain ») ; ce peut
être une abréviation de « Corollaire ».}

\page{9}
l'hom.\ naturel
\[
  \underline{\mathrm{Hom}}(P, \Omega)_s \otimes \mathcal{O}_{S,s}/(f_1, \ldots, f_n)
  \longrightarrow
  \underline{\mathrm{Hom}}_{\mathcal{O}_{S,s}}\bigl(P_s,\ \Omega_s/(f_1, \ldots, f_n)\Omega_s\bigr)
\]
soit un isom.

\emph{Dém.} On peut supposer $S$ affine, d'anneau $A$, donc
\uncertain{\ill{}} \struck{résolution} présentation
\[
  L_1 \to L_0 \to P
\]
avec $L_1$, $L_0$ libres de t.f., donc $\underline{\mathrm{Hom}}(P, Q)$ est le
$H^0$ du complexe
\[
  K^{\cdot} \colon\quad 0 \to \check{L}_0 \otimes Q \to \check{L}_1 \otimes Q \to 0
\]
(de degrés 0 et 1), tandis que $\mathrm{Hom}(P, Q \otimes M)$ est le $H^0$ du
complexe $K^{\cdot} \otimes M$. Une condition \uncertain{pour} que l'hom.\
naturel
\[
  H^0(K^{\cdot}) \otimes M \to H^0(K^{\cdot} \otimes M)
\]
soit bijectif est \uncertain{que l'on ait}
\[
  \mathrm{Tor}_1(M, H^1(K^{\cdot})) = \mathrm{Tor}_1(M, Z^1(K^{\cdot})) = 0 .
\]
Or $H^1(K) = R$ et $Z^1(K) = S$ sont des modules de t.f.\ sur $A$, auxquels on
peut appliquer le \struck{corollaire} lemme~6, b), en faisant
$M = \mathcal{O}_{S,s}/(f_1, \ldots, f_n)$, on gagne.
\note{la lettre notée ici $Q$ (dans $\underline{\mathrm{Hom}}(P, Q)$, dans
$K^{\cdot}$ et dans $Q \otimes M$) est un rond ouvert ; c'est sans doute le
$\Omega$ de l'énoncé, écrit vite. Dans $K^{\cdot}$ elle est lue sur ce qui
ressemble à un $P$. Les indications « degré 0 », « degré 1 » sont écrites
au-dessus des deux termes. Les lettres $R$ et $S$ (ce dernier sans rapport avec
le schéma $S$) sont lues telles qu'elles sont écrites.}

\page{10}
\struck{\emph{Lemme.} Soit $X$ un préschéma régulier} $\Longleftrightarrow$ les
$\Omega/(f_1, \ldots, f_n)\Omega$ sont \struck{\ill{} \ill{}} les injectifs
indécomposables \ill{}
\note{ligne d'en-tête biffée et surchargée ; au-dessus, en interligne :
« $\Omega$ \ill{} \ill{} \uncertain{définissables} par un C.M. ». Le passage
n'est lu qu'en partie.}

\emph{Lemme 8.} Soit $X$ un préschéma de Cohen--Macaulay, $\Omega$ un faisceau
fondamental sur $X$, alors pour toute suite $(f_1, \ldots, f_n)$ de sections
de $\mathcal{O}_X$ qui est $\mathcal{O}_X$-régulière,
$\Omega/(f_1, \ldots, f_n)\Omega$ est un faisceau fondamental sur le
préschéma \struck{noeth.}\ CM \struck{\ill{}} $V(f_1, \ldots, f_n)$.

\emph{\uncertain{Cor. Main} 9.} Soit $x \in X$ et $f_1, \ldots, f_n \in
\mathcal{O}_{X,x}$ un syst.\ de paramètres de $\mathcal{O}_{X,x}$, alors
$\Omega/(f_1, \ldots, f_n)\Omega$ est \struck{\ill{}} l'injectif
indécomposable dans la catégorie des modules sur
$\mathcal{O}_{X,x}/(f_1, \ldots, f_n)$.

\emph{\uncertain{Cor. Main} 10.} Soit $X$ un préschéma loc.\ noeth.\ qui est
localement immergeable dans un régulier. Alors il existe un ouvert dense $U$
de $X$, et un Module coh.\ $\Omega$ sur $X$, ayant les propriétés suivantes
\begin{itemize}
  \item[a)] Les anneaux locaux de $U$ sont C.M.
  \item[b)] Pour tout $x \in U$, et tout syst.\ de paramètres
    $f_1$, \ldots, $f_n$ de $\mathcal{O}_{X,x}$, le module
    $\Omega/(f_1, \ldots, f_n)\Omega$ sur
    $\mathcal{O}_{X,x}/(f_1, \ldots, f_n)$ est
\end{itemize}
\marginal{il suffit \uncertain{que} \ill{} soit excellent, ou seulement qu'il
satisfasse la condition de EGA IV 6.11.8.}
\note{la page est d'un papier plus foncé, mais la numérotation des énoncés
continue celle des pages précédentes. « localement immergeable dans un
régulier » est souligné d'un trait ondulé ; la note marginale, oblique dans la
marge de gauche, s'y rattache. Le « 10 » surcharge un autre numéro.}

\page{12}
l'injectif indécomposable.

Ces conditions impliquent la suivante
\begin{itemize}
  \item[c)] Pour tout homom.\ $u \colon R \to R'$ surjectif de modules de type
    fini sur un $\mathcal{O}_x$ ($x \in U$), tel que
    \[
      \mathrm{Hom}\bigl(R',\ \Omega_x/(f_1, \ldots, f_n)\Omega_x\bigr) \to
      \mathrm{Hom}\bigl(R,\ \Omega_x/(f_1, \ldots, f_n)\Omega_x\bigr)
    \]
    soit un isomorphisme pour tout syst.\ de paramètres $f_1, \ldots, f_n$ de
    $\mathcal{O}_{X,x}$, \uncertain{alors} $u$ est un isomorphisme.
\end{itemize}
\note{« pour tout syst.\ de paramètres \ldots\ de $\mathcal{O}_{X,x}$ » est
ajouté en interligne et renvoyé par un trait ; $R$ et $R'$ sont repassés.}

\struck{\emph{\uncertain{Cor. Main} 11.} \struck{Supposons} Supposons que
$U$, $\Omega$ soient comme dans Cor.~10. Soit $x \in U$, et soit
$u \colon R \to R'$ un homom.\ surjectif de modules de type fini sur
$\mathcal{O}_{X,x}$, \struck{\ill{} C.M.} tels que $R$ et $R'$ satisfassent la
condition suivante pour les foncteurs
$\mathfrak{L}_R(M) = \mathrm{Hom}(R, M)$ et
$\mathfrak{L}_{R'}(M) = \mathrm{Hom}(R', M)$ :}
\begin{itemize}
  \item[a)] \struck{$\mathfrak{L}(\Omega) \otimes A/(f_1, \ldots, f_n) \to
    \mathfrak{L}(\Omega \otimes A/(f_1, \ldots, f_n))$ est un isomorphisme
    pour les suites $A$-régulières $(f_1, \ldots, f_n)$ sur
    $\mathcal{O}_{X,x}$ (où $\mathfrak{L} = \mathfrak{L}_R$ ou
    $\mathfrak{L}_{R'}$) ;}
  \item[b)] \struck{$\mathfrak{L}_{R'}(\Omega) \xrightarrow{\ \sim\ }
    \mathfrak{L}_R(\Omega)$.}
\end{itemize}
\struck{Sous ces conditions, $u$ est un \emph{isom}.}
\marginal{à expliciter d'une autre façon}
\note{tout le Corollaire~11 est barré de longs traits obliques croisés. Dans
les définitions de $\mathfrak{L}_R$, $\mathfrak{L}_{R'}$, une première
écriture ($\mathrm{Hom}(R, \Omega \otimes M)$ ?) est biffée ; « suites
$A$-régulières » remplace « syst.\ de paramètres », biffé. La note marginale
est reliée par un trait vertical à la condition a).}

\page{14}
\struck{\emph{Cor.~11.} Supposons maintenant qu'on ait un syst.\ projectif
filtrant strict $(R_i)_{i \in I}$ de Modules coh.\ sur $U$, puis affine
$U = \operatorname{Spec} A$, et \struck{supposons} soit $F$ le foncteur sur
les $\mathcal{O}_U$-Modules $F = \ldots$}
\note{ce premier état est encadré et barré de longs traits obliques ; un
dernier « $F =$ » est laissé inachevé.}

\emph{\uncertain{Cor. Main} 11.} Soient $U$, $\Omega$ comme dessus,
\struck{\ill{} supposons} \uncertain{supposons} $S = U = \operatorname{Spec} A$,
soit $(P_i)$ un \struck{famille} syst.\ projectif de $A$-modules, et $F$ le
foncteur qu'il pro-représente. Supposons que l'on ait les deux conditions
suivantes
\begin{itemize}
  \item[(i)] $F(\Omega)$ est un $A$-Module de t.f.
  \item[(ii)] Pour tout $\mathfrak{p} \in S$ et tout syst.\ de paramètres
    $(f_1, \ldots, f_n)$ de $A_{\mathfrak{p}}$, l'application canonique
    \[
      F(\Omega) \otimes A_{\mathfrak{p}}/(f_1, \ldots, f_n) \to
      F\bigl(\Omega \otimes A_{\mathfrak{p}}/(f_1, \ldots, f_n)\bigr)
    \]
    est bijective.
\end{itemize}
Alors le syst.\ projectif $(P_i)$ est stationnaire.

De façon précise, (i) \struck{\ill{}} équivaut à dire que le système inductif
des $\mathrm{Hom}(P_i, \Omega)$ est stationnaire. Soit $i_0$ tel que
$\mathrm{Hom}(P_{i_0}, \Omega) \xrightarrow{\ \sim\ } \mathrm{Hom}(P_i, \Omega)$
pour $i \geq i_0$. Alors $P_i \simeq P_{i_0}$ pour $i \geq i_0$.

\emph{Dém.} \uncertain{Le} diagramme
\note{« le diagramme » est suivi d'un mot souligné en interligne et d'un
« $\varinjlim_i$ » renvoyé devant le premier terme ; dans la seconde colonne,
$\Omega/f\Omega_{\mathfrak{p}}$ est écrit
$\Omega_{\mathfrak{p}}/\underline{f}\,\Omega_{\mathfrak{p}}$ en bas.}

\begin{tikzcd}[column sep=small, nodes={font=\small}]
\varinjlim_i \mathrm{Hom}(P_i, \Omega) \otimes A_{\mathfrak{p}}/\underline{f} A_{\mathfrak{p}} \arrow[r, "\sim"] & \varinjlim_i \mathrm{Hom}(P_i, \Omega/\underline{f}\Omega_{\mathfrak{p}}) \\
\mathrm{Hom}(P_{i_0}, \Omega) \otimes A_{\mathfrak{p}}/\underline{f} A_{\mathfrak{p}} \arrow[r] \arrow[u] & \mathrm{Hom}(P_{i_0}, \Omega_{\mathfrak{p}}/\underline{f}\Omega_{\mathfrak{p}}) \arrow[u, "\alpha"]
\end{tikzcd}

\page{15}
montre que $\alpha$ est un isom., et comme les morphismes de transition de la
dernière $\varinjlim$ sont injectifs, il s'ensuit que
\[
  \mathrm{Hom}(P_{i_0}, \Omega/\underline{f}\Omega) \to
  \mathrm{Hom}(P_i, \Omega_{\mathfrak{p}}/\underline{f}\Omega_{\mathfrak{p}})
\]
est bijectif pour $i \geq i_0$, donc $P_i \to P_{i_0}$ localisé en
$\mathfrak{p}$ est un isom.\ (Cor.~10), donc $P_i \to P_{i_0}$ est bijectif.

Mettant ensemble prop.~3, lemme~5, \struck{lemme~6} Remarque~5 bis, et le
corollaire~10, on trouve
\note{« Remarque 5 bis » est écrit au-dessus de « lemme 6 », biffé.}

\emph{Théorème 12.} $A$ anneau noeth.\ quotient d'un régulier, $F$ un foncteur
$\mathrm{Mod}\,A \to (\mathrm{Ens})$. Pour que $F$ soit représentable par un
module de t.f., il faut et il suffit qu'il satisfasse aux conditions
suivantes :
\begin{itemize}
  \item[(ii)] $F$ exact à gauche ;
  \item[(i)] $F$ commute aux $\varinjlim$ filtrantes ;
  \item[(iii)] a) $\forall\, \mathfrak{p} \in \operatorname{Spec} A$, et tout
    $M$ de t.f.\ sur $\widehat{A}_{\mathfrak{p}}$,
    \[
      F(M) \xrightarrow{\ \sim\ } \varprojlim F(M/\mathfrak{m}_{\mathfrak{p}}^{n+1}M) ;
    \]
    b) $\forall\, f \in A$, et tout module de t.f.\ sur
    $\widehat{A} = \varprojlim A/f^{n+1}A$,
    \[
      F(M) \xrightarrow{\ \sim\ } \varprojlim F(M/f^{n+1}M) ;
    \]
  \item[(iv)] $\forall\, M$ de t.f.\ sur $A$, $F(M)$ est de t.f.\ sur $A$ ;
  \item[(v)] La fonction $\mathfrak{p} \mapsto
    \operatorname{rang}_{k(\mathfrak{p})} F(k(\mathfrak{p}))$ sur
    $\operatorname{Spec} A$ est constructible ;
  \item[(vi)] Pour tout quotient \ill{} $B$ de $A$, et tout $B$-Module de t.f.\
    $\Omega$, $\exists\, f \in B$, \uncertain{non} nilpotent, tel que pour
    tout $\mathfrak{p} \in \operatorname{Spec} B_f$, et tout système
    \struck{de} paramètres $f_1, \ldots, f_n$ de $B_{\mathfrak{p}}$,
    l'application canonique
\end{itemize}
\note{« par un module de t.f. » est ajouté sous la ligne et renvoyé par un
trait. Les deux premières conditions sont écrites dans l'ordre (ii), (i), le
numéro (i) de la seconde surchargeant un autre chiffre ; un trait courbe en
marge les relie, sans doute pour les intervertir. C'est l'exactitude à gauche
qui porte le numéro (ii), comme dans la Remarque~18 (page 21). En (vi), un mot
ajouté au-dessus de « quotient » n'est pas lu, et « nilpotent » est
surchargé ; la condition se poursuit page suivante.}

\page{16}
\[
  F(\Omega) \otimes B_{\mathfrak{p}}/\underline{f} B_{\mathfrak{p}}
  \longrightarrow
  F(\Omega \otimes B_{\mathfrak{p}}/\underline{f} B_{\mathfrak{p}})
\]
est bijective.

\emph{Dém.} La nécessité claire, celle de (vi) résultant du \struck{lemme} 6,
\uncertain{cor.\ Main} 7, et du fait qu'on peut choisir $f$ tel que $B_f$ soit
CM (alors $\underline{f} = (f_1, \ldots, f_n)$ est nécessairement une suite
\struck{\ill{}} $B_{\mathfrak{p}}$-régulière).

\emph{Suffisance.} On voit que les conditions envisagées sont stables par
passage à un $F_B$ ($B$ quotient de $A$). Cela nous permet de
\struck{supposer, par récurrence noeth., que $F_B$ est représentable par un
$B$-module de t.f., si $B \neq A$.} \struck{\ill{}} Grâce à prop.~3, on est
ramené à trouver un $f \in A$ non nilpotent, tel que $F_f$ soit repr.
\marginal{(i) (ii) (iii b) (iv)}

Notons que le lemme~5 et Remarque~5 bis impliquent déjà que $F$ est
pro-représentable par un système projectif strict $(P_i)$, les $P_i$ de t.f.
On \uncertain{peut} \uncertain{maintenant} \ill{} $f$, non nilpotent, tel que
le syst.\ $P_{i,f}$ soit stationnaire.
\marginal{(i) (ii) (iii a) (iv) (v)}

On conclut aussitôt par la \struck{condition} condition (vi), compte tenu du
cor.~11. \marginal{(vi)}

cqfd.
\note{les trois groupes de numéros de conditions sont écrits dans la marge de
gauche, en face des phrases qui les utilisent. « et Remarque 5 bis » et
« strict » sont ajoutés en interligne. Le passage biffé (deux lignes) est
encadré.}

\emph{Remarque 13.} Supposons que pour tout quotient intègre $B$ de $A$,
$\exists\, f \in B - 0$ tel que $B_f$ soit régulier. Cette condition rend
\struck{évidemment} inutile celle que $A$ soit quotient d'un régulier, et
entraîne que dans th.~12, (vi) $\Rightarrow$ (v), \struck{et dans (vi)}

\page{17}
\emph{Théorème 14.} Soient $f \colon X \to S$ propre de prés.\ finie, $F$, $G$
Modules de prés.\ finie sur $X$, avec $G$ $S$-plat. Alors,
\struck{\ill{}} il existe un Module $P$ sur $S$, de présentation finie, et un
isom.\ de foncteurs en le $\mathcal{O}_S$-Module quasi-cohérent $M$
\[
  \underline{\mathrm{Hom}}(P, M) \simeq
  f_{*}\,\underline{\mathrm{Hom}}(F,\ G \otimes_{\mathcal{O}_S} M)
\]
[En prenant les $\Gamma$, l'isom.\ fonctoriel en $M$
\[
  \mathrm{Hom}(P, M) \simeq \mathrm{Hom}(F,\ G \otimes_{\mathcal{O}_S} M),
\]
montre que $P$ est déterminé à isom.\ unique près comme représentant un
foncteur en $M$ des \uncertain{Modules} \ill{}.] La formation de $P$ commute
au changement de base quelconque $S' \to S$.
\marginal{\emph{NB}. Il y a le même \uncertain{énoncé} de
représentabilité, par un $W(\ldots)$ (d'où \ill{} \ill{}), pour $F$ de
prés.\ finie, $G$ quasi-cohérent ; \ill{} \ill{} \ill{} \ill{}
\uncertain{représentable} par un $P$ \ill{} \ill{} \ill{} de prés.\ finie
\ill{}.}
\note{la note marginale, oblique et serrée dans la marge de gauche, n'est lue
qu'en partie ; un « $F$ » y est repassé en surcharge. Le crochet ouvert devant
« En prenant les $\Gamma$ » est fermé après « près ».}

On voit facilement qu'on peut supposer $S = \operatorname{Spec} A$ affine, et
tout est une question de représentabilité d'un foncteur
$\mathrm{Mod}(A) \to (\mathrm{Ens})$. On est ramené comme d'habitude au cas
d'un $A$ de type fini sur $\mathbb{Z}$, donc quotient d'un régulier. On peut
donc faire appel au th.~12.

Les conditions (i) à (iv) sont standard. \struck{\ill{} \ill{}} En vertu de
la remarque~13, il reste à vérifier la condition (vi).
\struck{Condition (v) : \ill{} \ill{} \ill{} que $A$ est intègre,
$\exists$ ouvert non vide $U$ de $S = \operatorname{Spec} A$ tel que sur $U$,
la fonction $s \mapsto \dim \mathrm{Hom}(F_s, G_s)$ soit \uncertain{constante}.
Or \ill{} \ill{} $\mathrm{Hom}(F_s, G_s) = \Gamma(X_s,
\underline{\mathrm{Hom}}(F_s, G_s))$, d'autre part pour \ill{} On peut
supposer $B = A$.}
\note{« En vertu de la remarque 13, il reste à vérifier la condition (vi) »
est écrit en interligne au-dessus de deux lignes biffées ; tout le passage
suivant est encadré et barré de trois longs traits obliques, sauf « On peut
supposer $B = A$ », qui semble laissé.}

\page{18}
En fait, sauf pour la condition (ii) d'exactitude à gauche, nous pouvons
vérifier les conditions du th.~12 \struck{pour} pour un foncteur de la forme
\[
  M \longmapsto R^{q} f_{*}\bigl(\underline{\mathrm{Ext}}^{i}(F,\ G \otimes_{\mathcal{O}_S} M)\bigr)
\]
(sans supposer $G$ plat$/S$ ici) [ou aussi
$M \mapsto R^{q} f_{*}(R\,\underline{\mathrm{Hom}}(F,\ G \otimes_{\mathcal{O}_S} M))$],
i.e.\ dans le cas affine, pour des foncteurs
\[
  M \longmapsto H^{q}\bigl(X,\ \underline{\mathrm{Ext}}^{i}(F,\ G \otimes_{\mathcal{O}_S} M)\bigr)
\]
[ou $M \mapsto \mathrm{Ext}^{q}(X;\ F,\ G \otimes_{\mathcal{O}_S} M)$].
La vérification se fait en deux coups.
\note{« (sans supposer $G$ plat$/S$ ici) » est écrit en bout de ligne, à
droite ; le crochet fermant de la dernière formule est suivi d'un astérisque
sans renvoi visible.}

\emph{Lemme 15.} Soit $f \colon X \to S = \operatorname{Spec} A$, $f$ de t.f.\
et $A$ noeth., et $F$, $G$ \struck{\ill{}} des cohérents sur $X$, $\Omega$
cohérent sur $S$, $N$ entier. Alors il existe un ouvert dense $U$ de $S$ tel
que pour $\mathfrak{p} \in U$ et $f_1, \ldots, f_n$ \struck{\ill{}} suite
$A_{\mathfrak{p}}$-régulière, on ait
\[
  \underline{\mathrm{Ext}}^{i}(F,\ G \otimes_{\mathcal{O}_S} \Omega) \otimes A_{\mathfrak{p}}/\underline{f} A_{\mathfrak{p}}
  \xrightarrow{\ \sim\ }
  \underline{\mathrm{Ext}}^{i}(F,\ G \otimes_{\mathcal{O}_S} \Omega \otimes A_{\mathfrak{p}}/\underline{f} A_{\mathfrak{p}})
\]
pour $i \leq N$.
\note{l'énoncé est marqué d'un trait vertical dans la marge ; « entier » est
ajouté en interligne.}

\emph{Dém.} On peut supposer $X$ affine $= \operatorname{Spec} B$. Alors $F$ a
une résolution $L_{\cdot}$ par des modules libres de t.f.\ sur $B$, et
\[
  \underline{\mathrm{Ext}}^{i}(F,\ \underbrace{G \otimes_A \Omega}_{G'})
  = H^{i}\bigl(\underline{\mathrm{Hom}}^{\cdot}(L_{\cdot},\ G \otimes_{\mathcal{O}_S} \Omega)\bigr)
  = H^{i}\bigl(\underbrace{(\check{L}_{\cdot} \otimes_B G) \otimes_{\mathcal{O}_S} \Omega}_{K^{\cdot}}\bigr)
\]
tandis que
\[
  \underline{\mathrm{Ext}}^{i}(F,\ (G \otimes_{\mathcal{O}_S} \Omega) \otimes_A M) = H^{i}(K^{\cdot} \otimes M),
\]
et on veut des conditions pour que
\[
  \underline{\mathrm{Ext}}^{i}_{\mathcal{O}_X}(F, G') \otimes_{\mathcal{O}_S} M
  \xrightarrow{\ \sim\ }
  \underline{\mathrm{Ext}}^{i}_{\mathcal{O}_X}(F,\ G' \otimes_A M)
\]
\marginal{cette réduction \ill{} \ill{} \ill{} fait\ldots}
\note{dans la seconde formule, $\Omega$ est suivi d'un « $\otimes_A M$ »
surchargé ; l'accolade « $K^{\cdot}$ » est sous
$(\check{L}_{\cdot} \otimes G) \otimes \Omega$ ; les indices $A$ et
$\mathcal{O}_S$ sont souvent surchargés l'un de l'autre.}

\page{19}
i.e.
\[
  H^{i}(K^{\cdot}) \otimes_A M \xrightarrow{\ \sim\ } H^{i}(K \otimes_A M)
  \quad \text{si } i \leq N .
\]
Pour ceci, quitte à tronquer $K^{\cdot}$ en laissant tomber les $K^{j}$ pour
$j \geq N + 2$, on peut supposer $K^{\cdot}$ borné. Il suffit que l'on ait
\[
  \begin{cases}
    \mathrm{Tor}_i^{A}(K^{j}, M) = 0 & \text{pour } i > 0, \text{ tout } j \\
    \mathrm{Tor}_i^{A}(H^{j}(K^{\cdot}), M) = 0 & \text{pour } i > 0, \text{ tout } j.
  \end{cases}
\]
Comme les modules $K^{j}$, $H^{j}(K^{\cdot})$ sont de t.f., on est ramené à
\struck{prouver le} prouver le

\emph{Lemme 16.} [Soit $R$ cohérent sur $X$.] $\exists\, U$ dense dans
$S = \operatorname{Spec} A$ tel que $\mathfrak{p} \in U$, $f_1 \ldots f_n$
\struck{\ill{}} suite $A_{\mathfrak{p}}$-régulière, implique que
\[
  \mathrm{Tor}_i^{A}(F,\ A_{\mathfrak{p}}/\underline{f} A_{\mathfrak{p}}) = 0
  \quad \text{pour } i > 0 .
\]
\note{« Soit $R$ cohérent sur $X$. » est ajouté en interligne au-dessus de
l'énoncé, et relié par un trait ; la formule garde pourtant $F$.}

Or, pour prouver ceci, il suffit, d'après la démonstration du lemme~6, de
savoir que $\mathrm{gr}_J(\mathcal{O}_X)$ et $\mathrm{gr}_J(A)$ soient plats
sur $A/J$ aux points de $U$ ($J$ = nilradical de $A$). Or on peut choisir $U$
ainsi, par le \emph{th.\ de platitude générique}.
\note{« et $\mathrm{gr}_J(A)$ soient » est ajouté en interligne.}

Le lemme \struck{15} 15 montre que \struck{quitte à} \ill{}
\struck{\ill{} \ill{} \ill{} \ill{}, pour traiter le cas de}
$R^{q} f_{*}(\underline{\mathrm{Ext}}^{i}(F, G \otimes \Omega))$,
\struck{\ill{}} à prouver le
\note{la fin de la page est raturée et inachevée ; un « $\varinjlim$ » est
écrit au-dessus du mot biffé « quitte ».}

\page{20}
\emph{Lemme 17.} Soit $f \colon X \to S$ propre, $S$ noeth., et soit $E$
cohérent sur $X$. Alors il existe un ouvert dense $U$ de $S$ tel que
$\forall\, s \in U$, et $\underline{f} = f_1, \ldots, f_n$ suite
$\mathcal{O}_{S,s}$-régulière, l'hom.\ naturel
\[
  R^{j} f_{*}(E)_s \otimes \mathcal{O}_{S,s}/\underline{f}\,\mathcal{O}_{S,s}
  \longrightarrow
  R^{j} f^{(s)}_{*}\bigl(E^{(s)} \otimes \mathcal{O}_{S,s}/\underline{f}\,\mathcal{O}_{S,s}\bigr)
\]
soit un isom.\ pour tout $j$. ($f^{(s)} \colon X^{(s)} = X \times_S
\operatorname{Spec} \mathcal{O}_{S,s} \to \operatorname{Spec}
\mathcal{O}_{S,s}$, $E^{(s)}$ induit par $E$ sur $X^{(s)}$.)
\note{« $S$ noeth. » est ajouté en interligne ; l'énoncé est marqué d'un trait
vertical dans la marge.}

\emph{Dém.} Utilisant le lemme~16, on voit qu'on peut supposer déjà que
\[
  E^{(s)} \otimes \mathcal{O}_{S,s}/\underline{f}\,\mathcal{O}_{S,s}
  \simeq
  E^{(s)} \overset{\mathbb{L}}{\otimes} \mathcal{O}_{S,s}/\underline{f}\,\mathcal{O}_{S,s} .
\]
D'autre part, on peut supposer \struck{\ill{}} (faisant $X = S$, et
représentant $Rf_{*}(E)$ par un complexe borné à composantes cohérentes
$K^{\cdot}$ sur $S$) que
\[
  K^{\cdot} \otimes \mathcal{O}_{S,s}/\underline{f}\,\mathcal{O}_{S,s}
  = K^{\cdot} \overset{\mathbb{L}}{\otimes} \mathcal{O}_{S,s}/\underline{f}\,\mathcal{O}_{S,s} .
\]
Cela montre que l'hom.\ $\alpha$, qui est l'isom.\ de régularité
(où $M = A_{\mathfrak{p}}/\underline{f} A_{\mathfrak{p}}$)
\[
  Rf_{*}\bigl(\underbrace{E \overset{\mathbb{L}}{\otimes}_{\mathcal{O}_S} M}_{E \otimes_{\mathcal{O}_S} M}\bigr)
  \longrightarrow
  \underbrace{Rf_{*}(E) \overset{\mathbb{L}}{\otimes}_{\mathcal{O}_S} M}_{K^{\cdot} \overset{\mathbb{L}}{\otimes} M \,=\, K^{\cdot} \otimes M},
\]
s'identifie à
\[
  H^{j}(K^{\cdot}) \otimes M \to H^{j}(K^{\cdot} \otimes M) .
\]
Si on restreint encore $U$, on peut supposer
\[
  \mathrm{Tor}_i^{\mathcal{O}_S}(H^{j}(K^{\cdot}), M) = 0 \quad \text{pour } i > 0,\ \text{tout } j,
\]
et on sait que cela implique que l'hom.\ précédent est un isom.\ pour tout $j$.
cqfd.
\note{les produits tensoriels dérivés sont écrits $\otimes$ surmonté d'un
$\mathbb{L}$ ; « (où $M = A_{\mathfrak{p}}/\underline{f}A_{\mathfrak{p}}$) » est
ajouté en interligne. Dans les accolades, l'écriture est serrée et la lecture
de la seconde (« $K^{\cdot} \overset{\mathbb{L}}{\otimes} M = K^{\cdot}
\otimes M$ ») est incertaine.}

\end{document}
